\section{Boundary geometry and pointwise slit support}
\label{geom:section}

Boundary extension and uniqueness at a fresh tip identify the access
from a first disk hit. Pointwise smooth-slit support then supplies the
individual passages used later; uniformity over rough pasts will come
from the compactness argument in Section~\ref{rc:section}.

All closures at infinity in this section are spherical.  A filled hull
generated by a continuous curve means the curve together with the components
that it disconnects from the target.  The surviving domain is the component
towards the target.  We use the continuous-trace description of chordal
Loewner evolution; the claims below concern its boundary geometry, not the
existence of that trace. For the fixed parameter $0<\kappa<8$, we use
the continuous generating trace given by Rohde--Schramm
\cite[Theorem~5.1]{rohdeschramm2005}. Its domain Markov property follows
from composition of the Loewner maps and Brownian strong Markov.


\begin{lemma}[Continuous extension for a past domain]
\label{geom:extension}
Let $\eta:[0,T]\to\overline{\mathbb H}$ be continuous, with $\eta(0)=0$,
and let $D$ be the unbounded component of
$\mathbb H\setminus\eta[0,T]$.  Suppose $D$ is the surviving simply
connected domain of the trace.  Its conformal uniformization
$F:\mathbb H\to D$ extends continuously from the closed half-plane in
the spherical topology.  Boundary parameters of $F$ specify prime ends;
distinct parameters may have the same physical image.
\end{lemma}

\begin{proof}
Write
\[
 A=\bigl(\widehat{\mathbb C}\setminus\mathbb H\bigr)
       \cup\eta[0,T],
 \qquad K=\widehat{\mathbb C}\setminus D.
\]
The image of a compact interval under a continuous map is a locally
connected continuum.  This standard preservation statement is recalled in
\cite[Chapter~2, \S2.2, p.~19]{pommerenke1992}.
The other set in the union defining $A$ is a closed hemisphere, and the
two sets meet at $0$.  Their union is locally connected: at a point of
their intersection take small connected relative neighborhoods in each
and unite them; at a point belonging to only one, first choose a
neighborhood disjoint from the other closed set.  Thus $A$ is a locally
connected continuum.

The set $K$ is obtained from $A$ by adding complementary components.
Here is the elementary filling argument, stated locally in a Euclidean
chart.  Given a prescribed small diameter, uniform local connectedness of
the compact set $A$ supplies a scale $\delta$ at which nearby points of
$A$ can be joined by a connected subset of that diameter.  If $x,y$ are
points of the filled set with $|x-y|<\delta$ and
$\operatorname{dist}(x,A)>2\delta$, their joining segment avoids $A$
and remains in the same filled component as $x$.  Otherwise choose
nearest points $a,b\in A$ to $x,y$.  The segments $[x,a]$ and $[y,b]$
stay in their respective filled components until their endpoints, and
have lengths at most $2\delta$ and $3\delta$.  Also
$|a-b|\le6\delta$.  After reducing $\delta$, join $a$ to $b$ by a
small connected subset of $A$.  The resulting connected set joins
$x$ to $y$ with the prescribed small diameter.  A spherical version
follows by local charts, or by moving a point of $D$ to infinity.
This is also the filling statement in
\cite[\S2, Lemma~2.4, p.~7]{loridantluo2017}.

Consequently $K$ is locally connected.  The Torhorst boundary theorem,
in the form recorded in
\cite[\S2, Lemma~2.6, p.~8]{loridantluo2017}, implies that
$\partial D$ is locally connected.  The Continuity Theorem and Prime
End Theorem in \cite[Chapter~2, \S2.1, p.~18]{pommerenke1992} now give
the asserted continuous extension and its prime-end interpretation.
\end{proof}

\begin{lemma}[A fresh terminal point has one prime end]
\label{geom:fresh-tip}
In the preceding setting, suppose
$a=\eta(T)\in\mathbb H\cap\partial D$ and
$a\notin\eta[0,T)$.  Then $F$ has exactly one boundary parameter with
image $a$.  For the Loewner uniformization this parameter is the current
driving point.  Every access in $D$ tending to $a$ approaches that same
prime end.
\end{lemma}

\begin{proof}
First, $A\setminus\{a\}$ is connected: it is the union of the closed
lower half-plane with the connected image $\eta[0,T)$, and these meet
at $0$.  If $V$ is a filled complementary component of $A$, then
$\overline V\setminus\{a\}$ is connected.  Indeed the connected set
$V$ is dense in this set in the punctured sphere.  Moreover
$\partial V$ contains a point other than $a$.  A nonempty proper open
component with boundary contained in a singleton would be both open
and closed in the connected punctured sphere, and hence would equal
that entire punctured sphere, which is impossible here.  Thus
$\overline V\setminus\{a\}$ meets $A\setminus\{a\}$.  It follows
that
\[
 K\setminus\{a\}
  =(A\setminus\{a\})\,
    \cup\!\bigcup_{V\text{ filled}}
          (\overline V\setminus\{a\})
\]
is connected.  This argument permits arbitrary self-intersections of the
past away from $a$.

For completeness, the conformal consequence is a Jordan-curve
argument.  Replace the half-plane by the disk.  If two different circle
parameters $u,v$ had image $a$, join them by a simple crosscut in the
disk.  Its image is a Jordan curve $J$ contained in $D\cup\{a\}$:
the interior of the crosscut is mapped injectively into $D$, and its
only coinciding endpoints have image $a$.  Neither of the two open
circle arcs between $u$ and $v$ can map identically to $a$.  If an open
arc did, Schwarz reflection of $F-a$ across that arc would give an
analytic extension vanishing on an interval, contrary to univalence.
Choose a boundary image different from $a$ on each arc.  The two
crosscut components map to opposite sides of $J$, so those images lie
on opposite sides as well.  Both belong to $K\setminus\{a\}$,
whereas $J\cap K=\{a\}$.  This disconnects $K\setminus\{a\}$,
a contradiction.  This is the separation argument of
\cite[proof of Theorem~1, pp.~426--427, equations~(2.2)--(2.3)]{pommerenke1985},
applied to the complement rather than only to its boundary.

The continuous extension is onto $\overline D$, so there is a boundary
preimage of $a$ and it is unique.  Any sequence of domain points tending
to $a$ has inverse-image cluster points only in this singleton:
compactness of the closed disk and continuity of $F$ prove this claim.
This gives the assertion for every access.  The trace formula identifies
the current driving point as a preimage of the tip, so it is the unique
one.
\end{proof}

\begin{lemma}[Access at a first disk hit]
\label{geom:first-hit-access}
Let $\overline{B(z,r)}$ lie in the initial surviving domain, and write
$\eta$ for the future trace. Let $\tau$ be its first time reaching
$\partial B(z,r)$.
Suppose $\tau<\infty$ and $z$ belongs to the surviving domain $D_\tau$.
Then, writing $a=\eta(\tau)$, the segment $(a,z]$ is contained in
$D_\tau$ and approaches the current tip prime end at $a$.
If $g:D_\tau\to\mathbb H$ sends tip to $0$, target to infinity,
and is scaled so that $|g(z)|=1$, the image of $[a,z]$ is a continuous
path from $0$ to modulus $1$ whose inverse image is contained in
$\overline{B(z,r)}$.
\end{lemma}

\begin{proof}
The open ball is disjoint from the trace through time $\tau$.
It is connected and contains the surviving point $z$, so all of it
belongs to the same surviving component.  In particular $(a,z]\subset
D_\tau$.  The point $a$ has not previously been visited: an earlier
visit would have been an earlier hit of the circle.  It is also
disjoint from the initial hull by the hypothesis on the closed ball.
Lemma~\ref{geom:fresh-tip} therefore identifies the radial access
with the tip prime end, proving continuity of its image at $0$.
The other assertions follow from the normalization.  One may stop this
path at its first modulus-one point if desired.  No modulus of boundary
continuity, and no bound on its length, is needed.
\end{proof}

The same connected-component argument gives a useful convention about
killing: if a ball is disjoint from the trace when its center is
swallowed, its whole open interior is swallowed then.  Subsequent
trace growth in the surviving domain cannot enter that interior.

\begin{lemma}[Smooth slits have continuous drivers]
\label{geom:smooth-driver}
Let $\beta$ be a compact regular smooth simple arc, starting at $0$,
initially vertical, and otherwise contained in $\mathbb H$.  Its
successive prefixes, parametrized by half-plane capacity $2t$, are
generated by a continuous chordal Loewner driver on a compact interval.
\end{lemma}

\begin{proof}
Use initially any regular parameter $s$.  The complements of the
prefixes converge in the kernel sense when $s$ varies, so their
hydrodynamically normalized maps and inverse maps converge locally
uniformly.  Reflection near infinity, outside a ball containing the
whole slit, also gives convergence of their Laurent coefficients.

We record why the needed boundary convergence holds in this particular
family.  The complements consisting of the closed lower half-plane
and a prefix of $\beta$ are uniformly locally connected, in fact
uniformly locally arcwise connected.  On the simple compact arc,
uniform continuity of $\beta$ and of its inverse shows that nearby
points can be joined by a uniformly small subarc, which remains in
every prefix containing them.  Points of the arc sufficiently close
to the real line are on a uniformly short initial part of the arc,
since every closed later part has positive height.  The initial part
and the lower half-plane can thus be joined through $0$.  This proves
the asserted uniform local connectivity, including at the attachment.

Here is the standard small-crosscut proof that it implies the boundary
convergence required here.  Put the inverse maps on the disk using a
fixed half-plane-to-disk coordinate.  Their spherical area integrals
are bounded by the area of the sphere.  About any boundary parameter,
integrating over crosscuts formed by circles of radii between
$\delta$ and $\sqrt\delta$ gives a crosscut whose image has spherical
length at most $C/\sqrt{\log(1/\delta)}$, by Cauchy--Schwarz.  Its
endpoint images can be joined in the complement by an arc with
uniformly small diameter.  The resulting Jordan curve (or the image
crosscut alone when its endpoints coincide) cuts off a set of small
diameter on the side away from the image of an interior basepoint.
The latter images stay uniformly separated from the complements, by
kernel convergence and compactness of the parameter interval.
The disk cap at the boundary parameter maps into that small side.
This proves uniform equicontinuity at the boundary; interior
equicontinuity follows from local uniform convergence.  Hence the
inverse maps converge uniformly on the closed disk in the spherical
metric.  In particular their tip preimages converge: any cluster
point maps to the limiting tip, whose preimage is unique by
Lemma~\ref{geom:fresh-tip}.  Returning to hydrodynamic half-plane
coordinates shows that the tip coordinate $U(s)$ is continuous.
At $s=0$ the same conclusion follows from the initially vertical slit.

For clarity, the usual infinitesimal mapping argument now gives the
equation and the parameter.  If $g_s$ maps out a prefix and $s'>s$,
the inverse transition $h=g_s\circ g_{s'}^{-1}$ maps $\mathbb H$
into itself and has expansion $h(z)=z-m/z+O(|z|^{-2})$.  The
maximum principle gives $\operatorname{Im}(h(z)-z)\ge0$; its
half-plane harmonic representation, with reflection outside the
bounded incremental hull, is
\[
 h(z)=z-\int_{\mathbb R}\frac{d\nu(x)}{z-x},
 \qquad \nu\ge0,\qquad \nu(\mathbb R)=m.
\]
Here $m$ is the increment of half-plane capacity.  If $m=0$, the
representation gives $h(z)=z$, which is impossible for a nontrivial
added arc.  Capacity is therefore strictly increasing, and its
continuity follows from the Laurent-coefficient convergence above.

As $s'\downarrow s$, the added arc mapped by $g_s$ shrinks to $U(s)$.
Indeed every approach to the fresh tip has that same inverse-image
limit by Lemma~\ref{geom:fresh-tip}.  The support of $\nu$ consequently
shrinks to $U(s)$ as well.  One elementary way to see the last statement
is to enclose the added hull in a half-disk of radius $\epsilon$
about $U(s)$.  Its map-out function $Q=h^{-1}$ extends by reflection
to a normalized univalent function on $|z-U(s)|>\epsilon$.
Writing
\[
 Q(U(s)+\epsilon\zeta)
   =U(s)+\epsilon\left(\zeta+
                         \sum_{j\ge1}b_j\zeta^{-j}\right),
\]
Gronwall's exterior area theorem gives $\sum_{j\ge1}j|b_j|^2\le1$
\cite[Theorem~3.2.1 and Corollary~3.2.2, pp.~37--38]{beliaev2015}.
Cauchy--Schwarz at $\zeta=\pm2$ consequently gives
$|Q(U(s)\pm2\epsilon)-U(s)|<3\epsilon$.
The transition inverse is real analytic on the real rays beyond
these two images, so its representing measure is supported between
them.
Thus, uniformly on compact subsets of $\mathbb H$,
\[
 h(z)=z-\frac{m}{z-U(s)}+o(m).
\]
Invert this transition and use capacity $2t$ as parameter.  It follows
that
\[
 \frac{\partial g_t(z)}{\partial t}
       =\frac{2}{g_t(z)-U(t)}.
\]
The compact interval of capacity and continuity of $U$ prove the lemma.
\end{proof}

The following pointwise assertion is a special case of SLE support;
see Tran--Yuan \cite[Theorem~1.1 and Corollary~1.2]{tranyuan2020}.
We give a proof of the particular open-tube form needed here.

\begin{lemma}[Pointwise slit support]
\label{geom:slit-support}
Let $\beta$ be as in Lemma~\ref{geom:smooth-driver}, and let $q$ be its
terminal point.  Let $N$ be a relative open neighborhood of the whole
slit in $\overline{\mathbb H}$, including a relative neighborhood of
$0$, and let $V$ be an open neighborhood of $q$.  Chordal
$\mathrm{SLE}_{\kappa}$ satisfies
\[
 \mathbb P\{\text{for some }t<\infty,
       \ \gamma[0,t]\subset N\text{ and }\gamma(t)\in V\}>0.
\]
Suppose in addition that $G_1,\ldots,G_m$ are pairwise disjoint crosscuts
of $N$, each separating $0$ from all later crosscuts and from $V\cap N$.
On the displayed
event, the first hits of $G_1,\ldots,G_m$ occur in that order.
The probability asserted here is pointwise in the slit and tube; no
uniform positive lower bound over different slits or pasts is claimed.
\end{lemma}

\begin{proof}
Let $U$ be the deterministic slit driver on $[0,T]$.  We show that a
sufficiently small uniform neighborhood of $U$ forces the required
geometric event whenever the perturbed driver generates a continuous
trace.  Suppose $U_n\to U$ uniformly and let $K_n$ be their hulls at
time $T$.  The height bound $\operatorname{Im}K_n\le2\sqrt T$ and
the driver-extrema horizontal bound place all hulls in a common
bounded rectangle.

For each point of
$\overline{\mathbb H}\setminus\beta[0,T]$, its limiting forward
Loewner solution stays a positive distance from the driver up to
$T$.  This includes real points other than $0$: a simple slit
otherwise contained in $\mathbb H$ swallows no such point, and the
solution is reflected analytically across a surviving real interval.
Continuous dependence for the ODE supplies a whole surviving
neighborhood of each such point for all sufficiently large $n$.
Compactness of the fixed rectangle outside any prescribed neighborhood
of the slit now gives
\[
                  K_n\subset N\quad\hbox{for all large }n.
\]
The same argument applies to a smaller neighborhood when needed.

The reverse Loewner ODE gives local uniform convergence of the
inverse maps at time $T$: on compact subsets of $\mathbb H$ their
imaginary parts are bounded away from zero, so ordinary ODE
stability applies.  The corresponding domains consequently converge
in the kernel sense to $\mathbb H\setminus\beta[0,T]$.  For every
small $r>0$, this implies
\[
                 K_n\cap B(q,r)\ne\varnothing
                 \quad\hbox{for all large }n.
\]
Otherwise, along a subsequence the whole ball would survive.  A
compact passage from a fixed interior basepoint through the limiting
domain to a point of that ball, together with the ball, would then
belong to the kernel, incorrectly making $q$ an interior kernel
point.

This hull intersection also forces an intersection by the trace.
Choose $r$ with $\overline{B(q,2r)}\subset V\cap\mathbb H$, and
fix $w\in B(q,2r)\setminus\beta[0,T]$.  Upper hull containment
ensures $w\notin K_n$ eventually.  If the perturbed trace avoided
$B(q,2r)$, that connected ball would lie in one component of its
complement.  It meets $K_n$ by the preceding paragraph, so that
component would be swallowed.  The whole ball, including $w$, would
then belong to $K_n$, a contradiction.  Thus the trace visits $V$.
Since every earlier hull is contained in $K_n$, the trace stays in
$N$ up to that visit.  This sequential argument supplies an actual
positive driver tolerance by contradiction.

Brownian motion has positive probability in each uniform
neighborhood of the continuous function $U/\sqrt\kappa$, which
starts at zero.  This follows, for example, by polygonal approximation,
positive Gaussian increment probabilities, and positive Brownian-bridge
confinement probabilities on the finitely many resulting intervals.
Apply this to the driver $\sqrt\kappa B$ and use the almost-sure
continuous-trace property.  Finally, the stated crosscut separation
forces a continuous path in $N$ from $0$ to $V$ to hit the crosscuts
in order; backtracking within $N$ causes no problem.
\end{proof}

\begin{corollary}[Restarted support in open sets]
\label{geom:restart-support}
At a finite stopping time $\tau$ for ordinary chordal SLE in
$\mathbb H$, fix a realized past for which the conditional continuation
law $\mathbb P_\tau$ is defined.  Let $D$ be its surviving domain, and
let $F:\mathbb H\to D$ send $0$ to the current tip prime end and infinity
to the target.  Use its continuous spherical extension on
$\overline{\mathbb H}$.  Let
$\alpha:[0,1]\to\overline{\mathbb H}$ be a continuous simple path with
$\alpha(0)=0$ and $\alpha((0,1])\subset\mathbb H$.
If $O$ is relatively open in $\overline D$ and contains
$F(\alpha[0,1])$, and $V\subset D$ is an open neighborhood of
$F(\alpha(1))$, then the restarted half-plane trace $\widetilde\gamma$
satisfies
\[
 \mathbb P_\tau\{\text{for some }t<\infty,
   \ F(\widetilde\gamma[0,t])\subset O
       \text{ and }F(\widetilde\gamma(t))\in V\}>0.
\]
The same statement holds for a Euclidean similarity of the evolution.
The positive probability may depend on the realized past, $\alpha$,
$O$, and $V$.
\end{corollary}

\begin{proof}
Lemma~\ref{geom:extension} makes $N=F^{-1}(O)$ a relative open
neighborhood of $\alpha[0,1]$, including a relative half-disk at $0$.
The path $\alpha((0,1])$ lies in one component of $N\cap\mathbb H$,
and a short vertical segment from $0$ enters that component.  Polygonal
connectivity there joins this segment to a point of $F^{-1}(V)$.
Removing loops and smoothing corners inside $N$ gives a regular smooth
simple slit with an initially vertical segment and endpoint in
$F^{-1}(V)$.
Lemma~\ref{geom:slit-support}, applied under the restarted Brownian
Loewner law, gives the displayed event after mapping back by $F$.
For passages through successive physical fronts, one may choose $O$
as a thin tube with separating crosscuts and choose $V$ beyond the
last one; the final clause of that lemma then gives their order.
Any uniform lower bound over different pasts still requires a
separate argument.
\end{proof}
